\subsection{The rows and their cancellation}
The normalized average \eqref{eq:reconstruction-average} gives
\begin{equation}\label{eq:rg-reconstruction-definition}
 T_k(A)=b_k^{-1}\mathcal Q_k(A)
            \big(D_A+\mu+b_k\mathcal Q_k^\sharp(A)\mathcal Q_k(A)\big)
       =\mathcal Q_k(A)\big(b_k^{-1}D_A+1+\mu/b_k\big).
\end{equation}
The second form follows from
$\mathcal Q_k\mathcal Q_k^\sharp=1$.
At a connection where $F_\infty(A)$ is invertible, its mean map
$H_k(A)=b_kF_\infty(A)^{-1}\mathcal Q_k^\sharp(A)$ satisfies
\begin{equation}\label{eq:rg-reconstruction-left-inverse}
 T_k(A)H_k(A)
 =b_k^{-1}\mathcal Q_k(A)F_\infty(A)
               b_kF_\infty(A)^{-1}\mathcal Q_k^\sharp(A)=1.
\end{equation}
The barred identity follows where the inverse at $-A$ exists,
by conjugation with $\mathsf C$.

For $c\in\Delta_y$, let $c_j(c)$ be the center of its block
of side $L^j\eta_k$, $0\le j\le k$, so that $c_0(c)=c$ and
$c_k(c)=y$. The averaging phase is
\begin{equation}\label{eq:rg-reconstruction-phase}
 \varphi_A(y,c)=\exp\left(ie_k\eta_k
                 \sum_{j=0}^{k-1}(\tau A)(c_{j+1}(c),c_j(c))\right).
\end{equation}
The sum uses the average of the six rectilinear path sums.
We write $(\tau_kA)(y,c)$ for this sum.
All of these paths lie in $\Delta_y$ and their total length
in bonds is at most $\frac32(L^k-1)$. Thus
\begin{equation}\label{eq:rg-reconstruction-phase-moduli}
 \begin{aligned}
 |\varphi_A(y,c)^{\pm1}|&\le\exp(3\mathfrak i(A)/2),\\
 |u_A(c,c')^{\pm1}|&\le\exp(\eta_k\mathfrak i(A))
                                             \le\exp(\mathfrak i(A)).
 \end{aligned}
\end{equation}
for a photon bond, with $u_A(c,c')=\exp(ie_k\eta_k A(c,c'))$.

Write \(c_3=\pi x^{(3)}=c+\eta_kge_1\). The row form of
\eqref{eq:rg-reconstruction-definition} is
\begin{equation}\label{eq:rg-reconstruction-row}
 \begin{aligned}
 (T_k(A)\psi)(y)
 &=\int_{\Delta_y}\dd x\,\varphi_A(y,c)
                  \Bigl\{(1+\mu/b_k)\psi(x)\\
 &\qquad-(b_k\eta_k)^{-1}\sum_{i=0}^3M_i(x)
           [U_A(x,x^{(i)})\psi(x^{(i)})-\psi(x)]\Bigr\}.
 \end{aligned}
\end{equation}
Here $U_A=1$ for $i\le2$ and $U_A(x,x^{(3)})=u_A(c,c_3)$.
The normalized chamber integral over $\Delta_y$ has total mass
one. This row exists for every complex connection, regardless
of invertibility; the inverse hypothesis was needed only for
\eqref{eq:rg-reconstruction-left-inverse}.

The rows map antiperiodic fields to antiperiodic coarse fields:
translation by a period preserves blocks, neighbors and phases,
and supplies the single antiperiodic sign from the source field.

If $Y$ is a union of $M$-cubes and $y\in Y$ is a unit site,
then $\Delta_y\subset Y$: both $M$ and $L^k$ are odd, and the
unit block is contained in the $M$-cube of its center.
Define the two rows
\begin{equation}\label{eq:rg-reconstruction-two-pieces}
 \begin{split}
 (T_Y^{(1)}(A)\psi)(y)
 &=\int_{\Delta_y}\dd x\,\varphi_A(y,c)
    \Bigl\{(1+\mu/b_k)\psi(x)\\
 &\qquad-(b_k\eta_k)^{-1}\!\!\sum_{i:\,x^{(i)}\text{ over }Y}
           M_i(x)[U_A(x,x^{(i)})\psi(x^{(i)})-\psi(x)]\Bigr\},\\
 (T_Y^{(2)}(A)\chi)(y)
 &=-b_k^{-1}\eta_k^{\alpha-1}\int_{\Delta_y}\dd x\,
                1[c_3\notin Y]\varphi_A(y,c)M_3(x)\chi(x,x^{(3)}).
 \end{split}
\end{equation}
Rows outside $Y$ are zero. The first row contains ordinary
legs only over $Y$. The second row contains precisely the
type-$3$ pairs leaving $Y$, which are admissible boundary
pairs of length $\eta_k$. Their coefficients read $A$ only
on bonds with an endpoint in $\Delta_y$.
If $\chi=\delta_{\alpha,A}\psi$ on these boundary pairs, then
\begin{equation}\label{eq:rg-reconstruction-local-identity}
 T_Y^{(1)}(A)\psi+T_Y^{(2)}(A)\chi=T_k(A)\psi
                                             \quad\text{on the unit sites of }Y.
\end{equation}
This follows by replacing the missing type-$3$ difference in
\eqref{eq:rg-reconstruction-row} by
$\eta_k^\alpha\chi(x,x^{(3)})$.
The intra-cluster differences remain in $T_Y^{(1)}$.

\begin{proposition}[Uniform row bounds]\label{prop:rg-reconstruction-estimate}
The rows \eqref{eq:rg-reconstruction-two-pieces} and $T_k(A)$
obey the Euclidean norm estimates
\eqref{eq:rg-reconstruction-sup-estimate} for every complex
connection $A$. If $\mathfrak i(A),q_d(A)\le1$, their coefficient
row sums obey \eqref{eq:rg-reconstruction-coefficient-bound},
with the stated constant $C_T$ independent of $k,\ell,Y$.
\end{proposition}

\begin{proof}
Write $\mathfrak i=\mathfrak i(A)$ and $q_d=q_d(A)$.
Put $E=\Delta_y\times O_{\mathrm{cub}}$ and divide it into
\[
 E_1=\{x\in E:c_3\in\Delta_y\},\quad
 E_2=\{x\in E:c_3\in Y\setminus\Delta_y\},\quad
 E_3=\{x\in E:c_3\notin Y\}.
\]
For each coordinate direction $a\in\{\pm e_\mu\}$ there are
eight labels with $ge_1=a$. The face of $\Delta_y$ through
which $c+\eta_ka$ leaves the unit block has $\eta_k^{-2}$
photon sites. Thus
\begin{equation}\label{eq:rg-reconstruction-face-count}
 \#E=48\eta_k^{-3},\qquad \#(E_2\cup E_3)=48\eta_k^{-2}.
\end{equation}

The chamber hopping matrices satisfy
\[
 M_i(x^{(i)})-M_i(x)=w_i\gamma\cdot\widehat d_i(x),\qquad
 \sum_{i=0}^3w_i\gamma\cdot\widehat d_i(x)=0.
\]
Also $\|M_3(x)\|_2\le6$ and
$\|w_3\gamma\cdot\widehat d_3(x)\|_2=6$.
Reindex each intra-cluster sum by $x\mapsto x^{(i)}$ and
each type-$3$ sum in $E_1$ by $x\mapsto x^{(3)}$.
Both are involutions of their summation sets. The intra-cluster
sums become
$\sum_{x\in E}\varphi_A(y,c)\sum_{i\le2}
 (M_i(x^{(i)})-M_i(x))\psi(x)$, which equals
$-\sum_{x\in E}\varphi_A(y,c)w_3\gamma\cdot\widehat d_3(x)\psi(x)$.
The part over $E_1$ cancels with the corresponding matrix
difference from the reindexed type-$3$ sum. Thus
$(T_Y^{(1)}(A)\psi)(y)=R_0+R_1+R_2+R_3$, with
\begin{equation}\label{eq:rg-reconstruction-summation-parts}
 \begin{split}
 R_0&=(1+\mu/b_k)\frac{\eta_k^3}{48}
                         \sum_{x\in E}\varphi_A(y,c)\psi(x),\\
 R_1&=b_k^{-1}\frac{\eta_k^2}{48}
       \sum_{x\in E_2\cup E_3}
                     \varphi_A(y,c)w_3\gamma\cdot\widehat d_3(x)\psi(x),\\
 R_2&=-b_k^{-1}\frac{\eta_k^2}{48}\sum_{x\in E_2}
       \varphi_A(y,c)M_3(x)[u_A(c,c_3)\psi(x^{(3)})-\psi(x)],\\
 R_3&=-b_k^{-1}\frac{\eta_k^2}{48}\sum_{x\in E_1}
       [\varphi_A(y,c_3)u_A(c_3,c)-\varphi_A(y,c)]
                                             M_3(x^{(3)})\psi(x).
 \end{split}
\end{equation}
The face count and phase moduli bound the first three terms by
\[
 \exp(3\mathfrak i/2)\|\psi\|_{\mathrm{su}}
       \big[1+\mu/b_k+6b_k^{-1}
                           +6b_k^{-1}(1+\exp(\mathfrak i))\big].
\]
The fourth term requires the curvature estimate.

For neighboring photon sites $c,c'=c+\eta_ka$ in $\Delta_y$,
let $j(c,a)$ be the first scale at which their block centers
agree. The centers at all larger scales then agree as well.
For $j\ge2$, such a $c$ lies on a face of a block of side
$L^{j-1}\eta_k$. There are $L^{3(k-j+1)}$ such blocks inside
$\Delta_y$, with $L^{2(j-1)}$ sites on the indicated face.
For $j=1$ use all sites. We obtain
\begin{equation}\label{eq:rg-first-common-block-count}
 \#\{c:c,c+\eta_ka\in\Delta_y,\ j(c,a)=j\}
                         \le\eta_k^{-3}L^{-(j-1)}.
\end{equation}

Write
\[
 \begin{aligned}
 \varphi_A(y,c')u_A(c',c)-\varphi_A(y,c)
    &=\varphi_A(y,c)(\exp(v(c,a))-1),\\
 v(c,a)&=ie_k\eta_k
      [(\tau_kA)(y,c')+A(c',c)-(\tau_kA)(y,c)].
 \end{aligned}
\]
The contributions above scale $j(c,a)$ cancel in this bracket.
The remaining expression is the mean of the line sums over
closed paths from their common scale-$j$ center down to $c'$,
then across the bond to $c$, and back up the reversed paths.
All these loops lie in the scale-$j$ block, of side at most
$L^j\eta_k$, and have at most $3L^j$ bonds.
For a closed path of $n$ bonds in a box of side $s$, the axial
path estimate gives $|A(\Gamma)|\le2ns\|dA\|_\infty$.
To recall the reason, subtract the gauge obtained by line sums
from one corner in coordinate order $1,2,3$. The difference on
each bond is a sum over at most two strips of width one mesh
and length at most $s$. Its modulus is at most
$2s\|dA\|_\infty$, and the gauge sum around the loop vanishes.
This argument applies to complex fields as well.
Consequently, with $j=j(c,a)$,
\begin{equation}\label{eq:rg-reconstruction-phase-defect}
 |v(c,a)|\le6L^{2(j-k)}q_d,\qquad
 |\exp(v(c,a))-1|
                    \le6L^{2(j-k)}q_d\exp(6q_d).
\end{equation}
The factor $L^{2(j-k)}$ must be kept until after summing the
locations of the first common block. By
\eqref{eq:rg-first-common-block-count},
\begin{align*}
 \eta_k^2\sum_{\substack{c\in\Delta_y\\c+\eta_ka\in\Delta_y}}
                                  |\exp(v(c,a))-1|
 &\le6q_d\exp(6q_d)L^{1-k}\sum_{j=1}^kL^j\\
 &\le9Lq_d\exp(6q_d).
\end{align*}
The six directions, eight labels per direction, and
$\|M_3\|_2\le6$ now give
\[
 |R_3|_2\le54Lq_d\exp(6q_d)b_k^{-1}
                         \exp(3\mathfrak i/2)\|\psi\|_{\mathrm{su}}.
\]
Together with the first three terms this proves the first
line of \eqref{eq:rg-reconstruction-sup-estimate} for $T_Y^{(1)}$.
Take $Y$ to be the full torus to obtain the bound for $T_k$.

The second row is supported on $E_3$. Its definition and
\eqref{eq:rg-reconstruction-face-count} give directly
\[
 |(T_Y^{(2)}(A)\chi)(y)|_2
 \le b_k^{-1}\eta_k^{\alpha-1}\frac{\eta_k^3}{48}
            48\eta_k^{-2}\exp(3\mathfrak i/2)6\|\chi\|_{\mathrm{su}},
\]
which is the second bound.

For the coefficient norm, fix a row and choose each source
component to have the conjugate phase of its row coefficient.
Its component supremum is at most one, and its Euclidean
spinor supremum is at most two. The row coefficient sum is
then at most twice the corresponding bound in
\eqref{eq:rg-reconstruction-sup-estimate}. When
$\mathfrak i,q_d\le1$, use also $\mu\le1$ and $\eta_k^\alpha\le1$.
The resulting bound is $C_T$ as stated. The normalized index
measure is already in the row coefficients, so no extra
factor $48$ enters. The domain is invariant under sign, and
the matrix $\mathsf C$ is monomial; conjugating $T_Y^{(i)}(-A)$
or $T_k(-A)$ by it permutes row and column moduli.
This proves the barred bounds.
For a total connection in $\widetilde R_k+\mathfrak X_k$,
\eqref{eq:rg-total-imaginary-curvature-bounds} supplies
$\mathfrak i,q_d\le1$ under
\eqref{eq:rg-domain-charge-conditions}.
\end{proof}

The local reconstruction respects the symmetries of an activity.
Telescoping the scalar gauge sum in the averaging paths gives
\[
 \begin{aligned}
 \varphi_{A-d\lambda}(y,c)\exp(ie_k\lambda(c))
       &=\exp(ie_k\lambda(y))\varphi_A(y,c),\\
 u_{A-d\lambda}(c,c')\exp(ie_k\lambda(c'))
       &=\exp(ie_k\lambda(c))u_A(c,c').
 \end{aligned}
\]
Insert these two identities in
\eqref{eq:rg-reconstruction-two-pieces}. Every term acquires
the coarse factor $\exp(ie_k\lambda(y))$, where the boundary
leg transforms with the phase of its first point. Thus, with
$\lambda^{\mathrm u}$ the restriction to unit sites,
\begin{equation}\label{eq:rg-reconstruction-gauge-covariance}
 \begin{split}
 T_Y^{(1)}(A-d\lambda)\exp(ie_k\lambda)\psi
      &=\exp(ie_k\lambda^{\mathrm u})T_Y^{(1)}(A)\psi,\\
 T_Y^{(2)}(A-d\lambda)\chi^\lambda
      &=\exp(ie_k\lambda^{\mathrm u})T_Y^{(2)}(A)\chi,
 \qquad\chi^\lambda(x,x')=\exp(ie_k\lambda(c))\chi(x,x').
 \end{split}
\end{equation}
The identity also holds for gauges specified locally on the
bonds and sites read by the row.

A scheme symmetry $\tilde g=(S,a)$ maps nested blocks and their
centers to their images, preserves the order of each coordinate
path, and permutes the $48$ labels. Moreover
$M_i(\tilde gx)S=SM_i(x)$ and
$(\tilde gx)^{(i)}=\tilde g(x^{(i)})$.
Changing variables in the two rows therefore proves
\begin{equation}\label{eq:rg-reconstruction-lattice-covariance}
 T_{\tilde gY}^{(i)}(\tilde gA)\tilde g
                           =\tilde gT_Y^{(i)}(A).
\end{equation}
Complex conjugation changes the phases to those at $-A$, and
$\overline{M_i}=\varkappa^{-1}M_i\varkappa$, giving
\begin{equation}\label{eq:rg-reconstruction-reality}
 \overline{T_Y^{(i)}(\overline A)}
                         =\varkappa^{-1}T_Y^{(i)}(-A)\varkappa.
\end{equation}

The barred rows are $\bar T_Y^{(i)}(A)=\mathsf C^{-1}T_Y^{(i)}(-A)\mathsf C$.
Their gauges have the opposite sign. For the lattice covariance,
$(S^{-1})^T=\sigma(S)\mathsf C^{-1}S\mathsf C$ with
$\sigma(S)\in\{\pm1\}$: both matrix maps are homomorphisms,
and on each generating Dirac matrix they differ by a minus
sign. This same sign occurs in the source and target spin
actions and cancels in the barred intertwining identity.
For reality, $\mathsf C$ is real and
$\varkappa\mathsf C=-\mathsf C\varkappa$; the two signs from
commuting the conjugations cancel. Thus the barred versions
of \eqref{eq:rg-reconstruction-lattice-covariance} and
\eqref{eq:rg-reconstruction-reality} hold with the inverse
transpose spin action and the same $\varkappa$.


\subsection{Local substitution of the mean field}
\label{subsec:rg-step-local-substitution}

For $k\ge1$, let $Y$ be a union of partition $M$-cubes and allow independent
unit generators $\Psi,\bar\Psi$ over $Y$, together
with the ordinary and H\"older generators.
Using the two operators from
\eqref{eq:rg-reconstruction-two-pieces}, define
the substitution that fixes the latter generators and
sets, at the unit sites of $Y$,
\begin{equation}\label{eq:rg-step-local-substitution}
 \begin{split}
 \operatorname{sub}_{Y,B}\Psi
       &=T_Y^{(1)}(B)\psi+T_Y^{(2)}(B)\chi,\\
 \operatorname{sub}_{Y,B}\bar\Psi
       &=\bar T_Y^{(1)}(B)\bar\psi+\bar T_Y^{(2)}(B)\bar\chi.
 \end{split}
\end{equation}
The images are odd linear combinations of independent
generators and therefore anticommute with one another
and with the fixed generators. They determine a
unique Grassmann homomorphism.
Its new ordinary legs lie over $Y$ and its new pair
legs are boundary pairs of the second kind with first
point over $Y$. Its coefficients read $B$ on bonds
meeting $Y$, by the row support proved in
Section~\ref{subsec:rg-field-reconstruction}.
Since a bar index is preserved, the substitution
preserves balance and parity.

If an evaluation has
$\chi=\delta_{\alpha,B}\psi$ and the barred
relation on these boundary pairs, then
\eqref{eq:rg-reconstruction-local-identity}
turns \eqref{eq:rg-step-local-substitution}
into $T_k(B)\psi$ and $\bar T_k(B)\bar\psi$.
If in addition $\psi=H_k(B)\Phi$ and
$\bar\psi=\bar H_k(B)\bar\Phi$, the left-inverse
identity gives $\Phi$ and $\bar\Phi$ over $Y$.
Thus this local substitution represents the exact
mean field under evaluation.

The row covariances yield
\begin{equation}\label{eq:rg-step-substitution-covariances}
 \begin{split}
 a_g\operatorname{sub}_{gY,gB}
                        &=\operatorname{sub}_{Y,B}a_g,\\
 G_\lambda\operatorname{sub}_{Y,B^1}
                        &=\operatorname{sub}_{Y,B}G_\lambda,
                 \qquad B^1=B-d\lambda\text{ on bonds meeting }Y,\\
 \mathrm{Cg}\operatorname{sub}_{Y,-B}
                        &=\operatorname{sub}_{Y,B}\mathrm{Cg},\\
 \mathbf K\operatorname{sub}_{Y,-\overline B}
                        &=\operatorname{sub}_{Y,B}\mathbf K.
 \end{split}
\end{equation}
The final relation uses real $\mu$.
Both sides fix the old generators according to the
stated symmetry, so it suffices to compare their
images of $\Psi,\bar\Psi$. The lattice and real-gauge
relations follow from
\eqref{eq:rg-reconstruction-lattice-covariance} and
\eqref{eq:rg-reconstruction-gauge-covariance}; for
a local gauge equality the rows read only the
bonds on which that equality is assumed.
For charge conjugation use
$T_Y^{(i)}(-B)\mathsf C=\mathsf C\bar T_Y^{(i)}(B)$
and its barred counterpart. For the antilinear map,
\eqref{eq:rg-reconstruction-reality} with $-B$ gives
$\overline{T_Y^{(i)}(-\overline B)}
=\varkappa^{-1}T_Y^{(i)}(B)\varkappa$, and likewise
for the barred row. These identities give
\eqref{eq:rg-step-substitution-covariances} on
each new generator.

Assume $\mathfrak i(B),q_d(B)\le1$ under the standing hypotheses
of this section. Proposition~\ref{prop:rg-reconstruction-estimate}
gives the following norm statement. Assign weights
$w_1,w_2$ to the ordinary and pair legs and weight $w_3$
to the independent unit generators:
\begin{equation}\label{eq:rg-step-substitution-bound}
 \|\operatorname{sub}_{Y,B}F\|_{w_1,w_2}
                         \le\|F\|_{w_1,w_2;w_3}
             \quad\text{if }C_T(w_1+w_2)\le w_3.
\end{equation}
Indeed the image of a $\Psi$ generator has norm
at most $C_Tw_1+C_Tw_2$, and the original generators
retain their weights. The weighted linear-substitution
bound in Lemma~\ref{lem:linear-substitution}
applies on the subalgebra whose unit generators lie
over $Y$. The measures for the two row sums are
$\eta_k^3/48$ and $(\eta_k^3/48)^2$, respectively;
their identity kernels include the reciprocal
column measure. Thus this estimate has the same
normalization as the reconstruction bound.
